\documentclass[10pt]{article}
\usepackage[margin=0.75in]{geometry}
\usepackage{amsmath, amssymb, amsthm}
\usepackage{microtype}
\usepackage{enumitem}
\usepackage{fancyhdr}
\usepackage{titlesec}
\usepackage{tcolorbox}
\usepackage{booktabs}

\linespread{1.08}
\setlength{\parskip}{0.4ex plus 0.2ex minus 0.1ex}

\titleformat{\section}{\large\bfseries}{\thesection.}{0.5em}{}
\titleformat{\subsection}{\normalsize\bfseries}{\thesubsection}{0.5em}{}
\titlespacing*{\section}{0pt}{1.5ex}{1ex}
\titlespacing*{\subsection}{0pt}{1ex}{0.5ex}

\setlist{itemsep=1pt, topsep=3pt, parsep=1pt, leftmargin=1.5em}

\pagestyle{fancy}
\fancyhf{}
\fancyhead[L]{\small EECS 127/227AT}
\fancyhead[R]{\small Discussion 4}
\fancyfoot[C]{\small\thepage}
\renewcommand{\headrulewidth}{0.4pt}

\newcommand{\RR}{\mathbb{R}}
\newcommand{\NN}{\mathcal{N}}
\newcommand{\Range}{\mathcal{R}}
\DeclareMathOperator{\rank}{rank}
\DeclareMathOperator{\trace}{tr}

\begin{document}

\begin{center}
\begin{tabular}{lcr}
\textbf{EECS 127/227AT} & Optimization Models in Engineering & UC Berkeley \\
\textbf{Discussion 4} & & Spring 2026
\end{tabular}
\end{center}

\noindent\rule{\textwidth}{0.4pt}

\section{Gradient of the Cross Entropy Loss}

Consider the data $(\vec{x}_i, y_i)$ for $i = 1, \ldots, n$ where $\vec{x}_i \in \RR^d$ and $y_i \in \{0, 1\}$. Consider the parameter vector $\vec{w} \in \RR^d$. For each $i \in \{1, \ldots, n\}$, define the \emph{logistic function} $p_i \colon \RR^d \mapsto \RR$ given as
\begin{equation}
p_i(\vec{w}) = \frac{1}{1 + e^{-\vec{w}^\top \vec{x}_i}}.
\end{equation}

\begin{enumerate}[label=(\alph*)]
\item Find the gradient of the function $p_i(\vec{w})$.

\textbf{Solution:} The gradient is
\begin{equation}
\nabla p_i(\vec{w}) = \begin{bmatrix} \dfrac{\partial p_i}{\partial w_1}(\vec{w}) \\[6pt] \vdots \\[6pt] \dfrac{\partial p_i}{\partial w_d}(\vec{w}) \end{bmatrix}.
\end{equation}
Here
\begin{equation}
\frac{\partial p_i}{\partial w_j}(\vec{w}) = \frac{(\vec{x}_i)_j \, e^{-\vec{w}^\top \vec{x}_i}}{\left(1 + e^{-\vec{w}^\top \vec{x}_i}\right)^2}.
\end{equation}
Thus
\begin{equation}
\nabla p_i(\vec{w}) = \vec{x}_i \cdot \frac{e^{-\vec{w}^\top \vec{x}_i}}{\left(1 + e^{-\vec{w}^\top \vec{x}_i}\right)^2}
\end{equation}

\item For $i \in \{1, \ldots, n\}$, the \emph{cross entropy} of $p \in [0,1]$ against $y_i$ is defined as
\begin{equation}
H_i(p) \doteq -y_i \log(p) - (1 - y_i) \log(1 - p).
\end{equation}
Find the gradient of the function $\ell_i(\vec{w}) \doteq H_i(p_i(\vec{w}))$ with respect to $\vec{w}$.

\textbf{Solution:} The gradient is
\begin{equation}
\nabla_{\vec{w}} \ell_i(\vec{w}) = \begin{bmatrix} \dfrac{\partial \ell_i}{\partial w_1}(\vec{w}) \\[6pt] \vdots \\[6pt] \dfrac{\partial \ell_i}{\partial w_d}(\vec{w}) \end{bmatrix}.
\end{equation}
We can use the chain rule to find each component:
\begin{align}
\frac{\partial \ell_i}{\partial w_j}(\vec{w})
&= -\left[\frac{\partial H_i}{\partial p}(p_i(\vec{w}))\right] \left[\frac{\partial p_i}{\partial w_j}(\vec{w})\right] \\
&= -\left[\frac{y_i}{p_i(\vec{w})} - \frac{1 - y_i}{1 - p_i(\vec{w})}\right] \left[\frac{(\vec{x}_i)_j \, e^{-\vec{w}^\top \vec{x}_i}}{\left(1 + e^{-\vec{w}^\top \vec{x}_i}\right)^2}\right] \\
&= -\left[\frac{y_i}{1/(1 + e^{-\vec{w}^\top \vec{x}_i})} - \frac{1 - y_i}{e^{-\vec{w}^\top \vec{x}_i}/(1 + e^{-\vec{w}^\top \vec{x}_i})}\right] \left[\frac{(\vec{x}_i)_j \, e^{-\vec{w}^\top \vec{x}_i}}{\left(1 + e^{-\vec{w}^\top \vec{x}_i}\right)^2}\right] \\
&= -\left[y_i(1 + e^{-\vec{w}^\top \vec{x}_i}) - \frac{(1 - y_i)(1 + e^{-\vec{w}^\top \vec{x}_i})}{e^{-\vec{w}^\top \vec{x}_i}}\right] \left[\frac{(\vec{x}_i)_j \, e^{-\vec{w}^\top \vec{x}_i}}{\left(1 + e^{-\vec{w}^\top \vec{x}_i}\right)^2}\right] \\
&= -(\vec{x}_i)_j \left[y_i \frac{e^{-\vec{w}^\top \vec{x}_i}}{1 + e^{-\vec{w}^\top \vec{x}_i}} - (1 - y_i) \frac{1}{1 + e^{-\vec{w}^\top \vec{x}_i}}\right] \\
&= -(\vec{x}_i)_j \left[y_i(1 - p_i(\vec{w})) - (1 - y_i) p_i(\vec{w})\right] \\
&= -(\vec{x}_i)_j \left[y_i - p_i(\vec{w})\right] \\
&= (\vec{x}_i)_j \left[p_i(\vec{w}) - y_i\right].
\end{align}
Thus
\begin{equation}
\nabla_{\vec{w}} \ell_i(\vec{w}) = \vec{x}_i \left[p_i(\vec{w}) - y_i\right].
\end{equation}

\item Define the cross-entropy loss function as the sum of the cross entropy functions over the entire data set:
\begin{equation}
\ell(\vec{w}) = \sum_{i=1}^{n} \ell_i(\vec{w}).
\end{equation}
Find the gradient of the function $\ell(\vec{w})$.

\textbf{Solution:} Using linearity of the derivatives,
\begin{align}
\nabla \ell(\vec{w}) &= \sum_{i=1}^{n} \nabla \ell_i(\vec{w}) \\
&= \sum_{i=1}^{n} \vec{x}_i \cdot (p_i(\vec{w}) - y_i) \\
&= X^\top (\vec{p}(\vec{w}) - \vec{y}).
\end{align}
Here
\begin{equation}
X = \begin{bmatrix} \vec{x}_1^\top \\ \vdots \\ \vec{x}_n^\top \end{bmatrix}, \qquad
\vec{p}(\vec{w}) = \begin{bmatrix} p_1(\vec{w}) \\ \vdots \\ p_n(\vec{w}) \end{bmatrix}, \qquad
\vec{y} = \begin{bmatrix} y_1 \\ \vdots \\ y_n \end{bmatrix}
\end{equation}

Notice that this is the same type of gradient as least squares! All it requires is replacing our linear predictors $X\vec{w}$ with our logistic predictors $\vec{p}(\vec{w})$.
\end{enumerate}

\section{Gradients and Hessians}

The \emph{gradient} of a scalar-valued function $g \colon \RR^n \to \RR$, is the column vector of length $n$, denoted as $\nabla g$, containing the derivatives of components of $g$ with respect to the variables:
\begin{equation}
(\nabla g(\vec{x}))_i = \frac{\partial g}{\partial x_i}(\vec{x}), \quad i = 1, \ldots n.
\end{equation}

The \emph{Hessian} of a scalar-valued function $g \colon \RR^n \to \RR$, is the $n \times n$ matrix, denoted as $\nabla^2 g$, containing the second derivatives of components of $g$ with respect to the variables:
\begin{equation}
(\nabla^2 g(\vec{x}))_{ij} = \frac{\partial^2 g}{\partial x_i \, \partial x_j}(\vec{x}), \quad i = 1, \ldots, n, \quad j = 1, \ldots, n.
\end{equation}

For the remainder of the class, we will repeatedly have to take gradients and Hessians of functions we are trying to optimize. This exercise serves as a warm up for future problems. Compute the gradients and Hessians for the following functions:

\begin{enumerate}[label=(\alph*)]
\item Compute the gradient and Hessian (with respect to $\vec{x}$) for $g(\vec{x}) = \vec{y}^\top A \vec{x}$, where $A \in \RR^{m \times n}$ is a constant matrix and $\vec{y} \in \RR^m$ is a constant vector.

\textbf{Solution:} Let $A = \begin{bmatrix} \vec{a}_1 & \vec{a}_2 & \ldots & \vec{a}_n \end{bmatrix}$ where $\vec{a}_i$ is the $i$-th column of $A$. Then
\begin{align}
g(\vec{x}) &= \vec{y}^\top A \vec{x} \\
&= \vec{y}^\top \begin{bmatrix} \vec{a}_1 & \vec{a}_2 & \ldots & \vec{a}_n \end{bmatrix} \vec{x} \\
&= \vec{y}^\top (\vec{a}_1 x_1 + \vec{a}_2 x_2 + \ldots + \vec{a}_n x_n) \\
&= \sum_{i=1}^{n} (\vec{y}^\top \vec{a}_i) x_i.
\end{align}
Thus
\begin{equation}
\frac{\partial g}{\partial x_j}(\vec{x}) = \vec{y}^\top \vec{a}_j = \vec{a}_j^\top \vec{y},
\end{equation}
and the gradient $\nabla g(\vec{x}) = A^\top \vec{y}$. Since the gradient does not depend on $\vec{x}$, we then have the Hessian $\nabla^2 g(\vec{x}) = 0$.

\item Compute the gradient and Hessian of $h(\vec{x}) = \displaystyle\sum_{i=1}^{n} (x_i \log(x_i) - x_i)$ for $\vec{x} \in \RR^n_{++}$ and establish that the Hessian is positive semi-definite (as we will see soon in lecture, this establishes that $h$ is a convex function).

\emph{NOTE:} In fact, the Hessian is positive definite.

\textbf{Solution:} We have
\begin{align*}
\frac{\partial h(\vec{x})}{\partial x_i} &= \log(x_i) \\[4pt]
\frac{\partial^2 h(\vec{x})}{\partial x_i^2} &= 1/x_i \\[4pt]
\frac{\partial^2 h(\vec{x})}{\partial x_i \, \partial x_j} &= 0, \qquad \text{for } i \neq j.
\end{align*}
Hence the $i^{th}$ entry of $\nabla h(\vec{x})$ is $\log(x_i)$ and the Hessian $\nabla^2 h(\vec{x})$ is a diagonal matrix with the $(i,i)^{th}$ entry is $1/x_i$. As $x_i$ are positive, so is $1/x_i$ and so the diagonal matrix has only positive entries, and hence has positive eigenvalues.

\item Compute the gradient and Hessian of $g(\vec{x}) = e^{\vec{a}^\top \vec{x} + b}$ for $\vec{a}, \vec{x} \in \RR^n$, $b \in \RR$ and establish that the Hessian is positive semi-definite.

\textbf{Solution:} We can either compute the gradient and Hessian directly or we can use the properties of gradient and Hessians under composition with linear functions.

We will first see the former.
\begin{align*}
\frac{\partial g(\vec{x})}{\partial x_i} &= e^{\vec{a}^\top \vec{x} + b} a_i \\[4pt]
\frac{\partial^2 g(\vec{x})}{\partial x_i^2} &= e^{\vec{a}^\top \vec{x} + b} a_i^2 \\[4pt]
\frac{\partial^2 g(\vec{x})}{\partial x_i \, \partial x_j} &= e^{\vec{a}^\top \vec{x} + b} a_i a_j
\end{align*}
Writing these in matrix form, we get,
\begin{align}
\nabla g(\vec{x}) &= e^{\vec{a}^\top \vec{x} + b} \, \vec{a} \\
\nabla^2 g(\vec{x}) &= e^{\vec{a}^\top \vec{x} + b} \, \vec{a} \vec{a}^\top
\end{align}

The Hessian is clearly a rank one positive semi-definite matrix.

To see the second way, we notice that considering $e(x) = e^x$ for a scalar $x$, the derivative and second derivative of $e(x)$ is just $e^x$. Since the linear transform we are taking is $\vec{a}^\top \vec{x} + b$, we get the same result.
\end{enumerate}

\vfill
\begin{center}
\small \copyright\ UCB EECS 127/227AT, Spring 2026. All Rights Reserved. This may not be publicly shared without explicit permission.
\end{center}

\end{document}
